\documentclass[11pt,a4paper]{article}
\usepackage[T1]{fontenc}
\usepackage{lmodern,amsmath,amssymb,amsthm,mathtools,booktabs,microtype}
\usepackage[margin=26mm]{geometry}
\usepackage[colorlinks=true,linkcolor=blue,citecolor=blue,urlcolor=blue]{hyperref}
\setlength{\emergencystretch}{2em}
\newtheorem{theorem}{Theorem}[section]
\newtheorem{proposition}[theorem]{Proposition}
\newtheorem{lemma}[theorem]{Lemma}
\newtheorem{corollary}[theorem]{Corollary}
\theoremstyle{definition}\newtheorem{definition}[theorem]{Definition}
\theoremstyle{remark}\newtheorem{remark}[theorem]{Remark}
\newcommand{\Q}{\mathbb Q}\newcommand{\Z}{\mathbb Z}
\newcommand{\R}{\mathbb R}\newcommand{\A}{\mathbb A}
\newcommand{\PP}{\mathbb P}\newcommand{\sqcl}[1]{[#1]}
\DeclareMathOperator{\Tr}{Tr}\DeclareMathOperator{\Nm}{N}
\DeclareMathOperator{\N}{N}\DeclareMathOperator{\rad}{rad}
\DeclareMathOperator{\sfpart}{sf}
\title{A Valuation Obstruction to Products of Elliptic Abscissas}
\author{}\date{16 November 2026}
\begin{document}\maketitle
\begin{abstract}
If the constant coefficient of a local integral short Weierstrass equation
is a nonsquare unit at an odd prime $p$, every rational abscissa has
even $p$-adic valuation. The same restriction holds for every finite
product of integral powers of such abscissas. In particular,
$y^2=x^3-2$ has positive rational rank, but $5$ does not belong to
the multiplicative group generated by its rational abscissas.
The proof uses an elementary valuation argument and finite-field
torsion counts.
\end{abstract}

\section{Scope}
For an elliptic curve $E/\Q$, let $X(E)$ be the set of its nonzero
rational affine abscissas. We study the subgroup
\[
 \langle X(E)\rangle=
 \left\{\prod_{i=1}^r x_i^{n_i}:
       r\ge0,\ x_i\in X(E),\ n_i\in\Z\right\}\subseteq\Q^\times.
\]
The number of factors is not bounded. The result below supplies an
element outside this entire subgroup for an explicit positive-rank
curve. It concerns multiplicative combinations and imposes no
restriction on sums or on more general existential constructions.
Obstructions for fixed rational functions of finitely many elliptic
points are treated separately in~\cite{EllipticRatios}.

\section{A valuation obstruction}
\begin{lemma}\label{lem:abscissa-parity}
Let
\[
                       E:y^2=x^3+Ax+B
\]
be a nonsingular elliptic curve over $\Q$. Suppose that $p$ is
an odd prime, $A$ and $B$ are $p$-integral, and the reduction
of $B$ is a nonsquare in $\mathbb F_p^\times$.
Every affine rational point of $E$ has nonzero abscissa $x$
and even valuation $v_p(x)$. The multiplicative subgroup of
$\Q^\times$ generated by all these abscissas consists entirely
of numbers of even $p$-adic valuation.
\end{lemma}
\begin{proof}
If $x=0$, then $y^2=B$, contrary to the nonsquare residue of
$B$. If $v_p(x)>0$, the equation has a unit right-hand side
congruent to $B$ modulo $p$, giving the same contradiction.
If $v_p(x)<0$, then
\[
        3v_p(x)<v_p(Ax),\qquad 3v_p(x)<v_p(B).
\]
The term $x^3$ has the unique least valuation on the right,
so $2v_p(y)=3v_p(x)$. This forces $v_p(x)$ to be even.
The remaining case has valuation zero. Valuations add under
multiplication and change sign under inversion, proving the
last assertion.
\end{proof}

\begin{proposition}\label{prop:abscissa-example}
The curve $E:y^2=x^3-2$ has a rational point of infinite order,
but $5$ is not a ratio of two rational abscissas of $E$, or
any finite product of integral powers of its rational abscissas.
\end{proposition}
\begin{proof}
The point $(3,5)$ lies on $E$. The curve has good reduction at
$5$ and $7$, since its discriminant is $-1728$. Direct counting
gives
\[
                       \#E(\mathbb F_5)=6,
                       \#E(\mathbb F_7)=7.
\]
For completeness, modulo $5$ the numbers of ordinates above
$x=0,1,2,3,4$ are $0,2,2,1,0$. Modulo $7$ the corresponding
numbers are $0,0,0,2,0,2,2$. Add the identity in each case.

Prime-to-$p$ rational torsion injects under good reduction at
$p$; see also~\cite[Chapter VII, Section 3]{Silverman}. One way to see the needed injectivity is to use the formal
parameter $z=-x/y$ in the reduction kernel: if $p\nmid m$,
then $[m]z=mz+O(z^2)$ has the same finite valuation as $z$ and
cannot vanish. Thus a prime divisor $\ell\ne5,7$ of a rational
torsion order would divide both $6$ and $7$, which is impossible.
The $5$-primary torsion is excluded by reduction at $7$, and
the $7$-primary torsion by reduction at $5$. Hence $E(\Q)$
has trivial torsion, and $(3,5)\ne O$ has infinite order.

At $p=5$, the constant coefficient is $B=-2\equiv3$, a
nonsquare. Lemma~\ref{lem:abscissa-parity} shows that every
multiplicative combination in the statement has even
$5$-adic valuation, whereas $v_5(5)=1$.
\end{proof}

\begin{thebibliography}{99}
\bibitem{Silverman}
J.~H.~Silverman, \emph{The Arithmetic of Elliptic Curves},
second edition, Graduate Texts in Mathematics 106,
Springer, 2009.

\bibitem{EllipticRatios}
\emph{Locally Dense Diophantine Sets in the Rational Subfield},
project manuscript, 9 November 2026,
\nolinkurl{negative_results/papers/number_fields/elliptic_ratios_and_integrality_20261109_v13.tex}.
\end{thebibliography}
\end{document}
